\chapter{Einleitung}
\label{cap:einleitung}



\section{Problemstellung}

Beispiel für \cite{Akgun07}.. 

\lipsum[2-4]

\section{Ziel der Arbeit}
\lipsum[2-4]




\section{Aufbau}
Die Arbeit ist in Kapitel unterteilt, welche so weit als möglich in sich selbst abgeschlossen sind. 

\begin{description}
\item[Kapitel \ref{cap:einleitung}] gibt einen Überblick über die Problemstellung und das Ziel der Arbeit. 
\item[Kapitel \ref{cap:theorie}] führt den Leser in das Themengebiet der Softwareentwicklung und erlernen von Sprachen ein. Das Wissen über verschiedene Verfahren zum erstellen von Feedback in Sprachen soll aufbauend vermittelt werden.
\item[Kapitel \ref{cap:empirisch}] befasst sich mit der konkreten Umsetzung eines/oder mehrerer Unterrichtskonzeptes aufbauend auf den vorher elaborierten Theorieteil. 
\item[Kapitel \ref{cap:conclusion}] dient als eine Zusammenfassung der Arbeit und gibt Ausblick auf mögliche weitere Forschungsfelder.
\end{description}